\documentclass[11pt,a4paper]{article}
\usepackage[margin=28mm]{geometry}
\usepackage{amsmath,amssymb,amsthm,booktabs,array}
\usepackage[T1]{fontenc}
\usepackage{lmodern}
\usepackage{microtype}
\usepackage{xurl}
\usepackage{xcolor}
\usepackage[colorlinks=true,linkcolor=blue!50!black,urlcolor=blue!50!black,citecolor=blue!50!black]{hyperref}
\hypersetup{pdftitle={A computer-assisted sharp diagonal Hlawka inequality for p at least 21},pdfauthor={Yuxuan Xu}}
\newtheorem{theorem}{Theorem}[section]
\newtheorem{lemma}[theorem]{Lemma}
\newtheorem{proposition}[theorem]{Proposition}
\newtheorem{corollary}[theorem]{Corollary}
\theoremstyle{remark}\newtheorem{remark}[theorem]{Remark}
\newcommand{\R}{\mathbb R}\newcommand{\C}{\mathbb C}
\newcommand{\Np}{N_p}\newcommand{\F}{F_{p,K}}
\newcommand{\repo}{\url{https://github.com/Anchen-AI4Math/hlawka-diagonal-certificates}}
\DeclareMathOperator{\Rea}{Re}
\allowdisplaybreaks
\setlength{\emergencystretch}{2em}
\title{A computer-assisted sharp diagonal Hlawka inequality\protect\\for real exponents $p\ge21$}
\author{Yuxuan Xu\thanks{AI-assisted research and manuscript preparation; see the disclosure at the end.}\\
School of Computer Science, Wuhan University}
\date{17 October 2026\\\small Research preprint, version 1}
\begin{document}
\maketitle
\begin{abstract}
We give a computer-assisted proof that, for every real $p\ge21$, the
dimension-independent sharp Hlawka constant of complex $\ell_p$ is the
maximum of an explicit one-parameter ratio on the cyclic triple
$(-t,1,1),(1,-t,1),(1,1,-t)$, $1/2\le t\le2$.
The new finite calculation covers $21\le p\le32$; the range $p\ge32$
uses a previously accepted Prove2Me theorem. The improvement uses norm
caps, clipped pair-coordinate tails, and linear coordinate-sum constraints
on a convex averaging domain, together with a sharper residual Hessian
estimate. Rational linear-programming dual witnesses and interval
certificates cover real exponent intervals, rather than sampled exponents.
The finite obligations pass both Arb and a standard-library dyadic rational
implementation. We give the analytic reduction, certificate specification,
and reproduction instructions. This preprint has not undergone independent
analytic review; the new range is not Lean-formalized, and no optimality of
the exponent threshold $21$ is asserted.
\end{abstract}
\noindent\textbf{Keywords:} Hlawka inequality; $\ell_p$ norm; sharp constant;
computer-assisted proof; interval arithmetic; convexity.

\section{Introduction and statement}
For $p>1$ and $v\in\C^n$ write $\Np(v)=(\sum_i|v_i|^p)^{1/p}$.
For a triple $(x,y,z)$ put
\begin{align}
S&=\Np(x)+\Np(y)+\Np(z),& T&=\Np(x+y+z),\nonumber\\
P&=\Np(x+y)+\Np(x+z)+\Np(y+z),& G&=2S-P.\label{eq:defs}
\end{align}
The Hlawka constant problem asks for the least $K$, uniformly in the
dimension, such that $S-T\le KG$. Equivalently, the deficit
\begin{equation}\label{eq:F}
\F=(2K-1)S+T-KP=KG-(S-T)
\end{equation}
must be nonnegative. The triangle inequality gives $G\ge0$.
For diagonal complex matrices this is exactly the corresponding Schatten
$p$-norm question, since their singular values are the absolute values of
their diagonal entries. We do not determine the sharp constant for general
noncommuting matrices.

Audenaert and Kittaneh formulate a dimension-independent operator version
in Section 8.2, Problem 7 of \cite{AK}. A survey of the classical inequality
and related questions is \cite{MN}. The diagonal construction and its
reduction to a cyclic family were developed in the formalization project
\cite{savarin}; the Prove2Me campaign records an accepted exponent cutoff
$32$ \cite{campaign,solution32}. We use that result for the unbounded tail.
Our contribution is an explicit certified extension to $[21,32]$, retaining
additional convex constraints in the localization argument. We reproduce
the analytic reductions needed for this interval below rather than requiring
the reader to combine the earlier working notes.

Define, for $1/2\le t\le2$,
\begin{align}
A_p(t)&=(t^p+2)^{1/p},&B_p(t)&=(2|1-t|^p+2^p)^{1/p},\nonumber\\
R_p(t)&=\frac{3A_p(t)-3^{1/p}|2-t|}{6A_p(t)-3B_p(t)},&
K_p&=\max_{1/2\le t\le2}R_p(t).\label{eq:cyclic}
\end{align}
Each pair of vectors in the cyclic family is nonproportional. Strict
convexity of $\ell_p$ therefore makes the denominator positive. Thus the
maximum exists by continuity and compactness. Since $R_p(2)=1$, $K_p\ge1$.

\begin{theorem}\label{thm:main}
For every real $p\ge21$, every $n\ge1$, and every $x,y,z\in\C^n$,
\begin{equation}\label{eq:main} S-T\le K_p(2S-P).\end{equation}
The constant $K_p$ is the least constant valid uniformly in $n$.
For each fixed $n\ge3$, it is also the least constant in that dimension.
\end{theorem}
The proof on $[21,32]$ consists of the analytic arguments below and the
finite certificate lemma in Section~\ref{sec:cert}. The proof on $[32,\infty)$
invokes \cite{solution32}, whose downloaded source and acceptance metadata
are preserved in the repository. That source was not independently rebuilt
in this project. The theorem is a computer-assisted claim in this preprint,
not a claim that the platform has accepted the cutoff $21$.

\section{Dimension reduction and the scalar envelope}\label{sec:reduce}
We first explain why it suffices to exclude real three-dimensional
counterexamples. Throughout, $p>1$ and $K\ge1$.

\begin{lemma}\label{lem:dim}
If $\F\ge0$ for all triples in $\R^3$, it holds in every finite real
dimension and then in every finite complex dimension.
\end{lemma}
\begin{proof}
For a real triple remove coordinates in which all three entries vanish.
Give each remaining coordinate a nonnegative weight $w_i$, preserving the
three pair moments
\[
\sum_i w_i\bigl(|x_i+y_i|^p,|x_i+z_i|^p,|y_i+z_i|^p\bigr)
=\sum_i\bigl(|x_i+y_i|^p,|x_i+z_i|^p,|y_i+z_i|^p\bigr).
\]
The negative part $-KP$ is constant, and the positive part of $\F$ is
concave in $w$, because $s^{1/p}$ is concave and $2K-1\ge0$.
If more than three weights are positive, the corresponding moment columns
are linearly dependent. Each column is nonnegative and nonzero: three zero
pair entries would imply $x_i=y_i=z_i=0$. Any nonzero kernel direction
therefore has both signs. Its feasible interval has two finite endpoints,
each with fewer positive weights. Concavity implies that at least one
endpoint does not increase $\F$. Starting from all weights equal to one
and iterating leaves at most three positive weights. Multiply coordinate
$i$ by $w_i^{1/p}$ and delete zero coordinates. A counterexample would thus
give a counterexample in $\R^3$, padding with zeros when necessary.

For complex vectors put $a_p=(2\pi)^{-1}\int_0^{2\pi}|\cos\theta|^p,d\theta>0$.
For positive integers $M$, define the real-linear map
\[
(Q_Mx)_{m,i}=M^{-1/p}\Rea(e^{2\pi im/M}x_i),\qquad 0\le m<M.
\]
Riemann-sum convergence and rotation invariance give
$\Np(Q_Mv)\longrightarrow a_p^{1/p}\Np(v)$ for each of the seven vectors
in \eqref{eq:defs}. Apply the real inequality in dimension $Mn$, take
the limit, and divide by $a_p^{1/p}$. This uses convergence of norms,
not interchange of a norm with an integral.
\end{proof}

\begin{lemma}[Weighted three-point convexity]\label{lem:weighted}
Let $f:\R\to\R$ be convex, let $a_i>0$, $s=\sum_{i=1}^3a_i$, and
$m=s^{-1}\sum_i a_i u_i$. Then
\[
\sum_{i<j}(a_i+a_j)f\left(\frac{a_iu_i+a_ju_j}{a_i+a_j}\right)
\le\sum_i a_if(u_i)+s f(m).
\]
\end{lemma}
\begin{proof}
Order $u_1\le u_2\le u_3$; the case of three equal points is immediate.
If $u_2\le m$, the means for pairs $(1,3)$ and $(2,3)$ lie in $[m,u_3]$.
Their total mass is $s+a_3$ and their first moment is $sm+a_3u_3$.
The chord bound on this interval bounds their combined contribution by
$sf(m)+a_3f(u_3)$. Jensen bounds the remaining pair by
$a_1f(u_1)+a_2f(u_2)$. If $m\le u_2$, use the chord on $[u_1,m]$
for pairs $(1,2)$ and $(1,3)$ and Jensen for $(2,3)$.
\end{proof}

Set
\begin{equation}\label{eq:envelope}
g_p(u)=((1+u^p)/2)^{1/p},\qquad
E_p(u)=\frac{1-u}{2(1-g_p(u))}\quad(0\le u<1).
\end{equation}
For a real triple normalized by $S=1$ and having no zero vector,
Lemma~\ref{lem:weighted}, applied coordinatewise to $f(s)=|s|^p$
with weights $\Np(x),\Np(y),\Np(z)$, gives
\[
\sum_{i<j}\frac{\Np(x_i+x_j)^p}{(\Np(x_i)+\Np(x_j))^{p-1}}\le1+T^p.
\]
Jensen with weights $(\Np(x_i)+\Np(x_j))/2$ implies $P\le2g_p(T)$.
Here $x_1,x_2,x_3$ denote the three vectors, not scalar coordinates.
The function $g_p$ is convex and $g_p(1)=1$, so its secant slope
$(1-g_p(u))/(1-u)$ is nondecreasing; $E_p$ is nonincreasing.
A zero vector gives $\F=(K-1)$ times the remaining pair gap and cannot
be a counterexample. For a counterexample $T<1$, and
\begin{equation}\label{eq:envelope-use}
K<\frac{1-T}{G}\le E_p(T).
\end{equation}
Consequently, $E_p(q)<K_p$ forces $T<q$ when $S=1$.

\section{Normalization, orientation, and coarse localization}\label{sec:coarse}
Fix one of the two windows in Table~\ref{tab:profiles}; write its endpoints
as $p_0,p_1$, and put $A=3q$, $k_0=cp_0$. All terminating decimals
in this paper and the input profiles denote exact rationals.
\begin{table}[ht]
\centering\small
\begin{tabular}{lcc}\toprule
Parameter & Window I & Window II\\\midrule
$[p_0,p_1]$ & $[21,26]$ & $[26,32]$\\
$c$ & $273/625$ & $443/1000$\\
$q$ & $1011/2500$ & $197/500$\\
Orientation error $e$ & $13/100$ & $1/10$\\
Initial tail error $\varepsilon_0$ & $10^{-5}$ & $10^{-5}$\\
Refined tail error $\varepsilon$ & $10^{-8}$ & $10^{-8}$\\
Coarse $(\ell_0,u_0,v_0,w_0)$ & $(.378,1.214,.698,1.214)$ & $(.477,1.183,.750,1.183)$\\
Young parameters $(e_1,e_2)$ & $(1,1/4)$ & $(1/10,1/10)$\\
Slab width & $1/32$ & $1/16$\\\bottomrule
\end{tabular}
\caption{Exact profiles used by the finite verification.}\label{tab:profiles}
\end{table}
The certificates establish
\begin{equation}\label{eq:Kbounds} cp<K_p\le p/2,\qquad E_p(q)<K_p.\end{equation}
Their verification is specified in Section~\ref{sec:cert}.

The reflection $(x,y,z)\mapsto(-(x+y+z),y,z)$ exchanges $T$ and
$\Np(x)$, permutes the pair norms, and changes $\F$ by
$(2K-2)(T-\Np(x))$. Reflect in a largest singleton norm if necessary.
A counterexample can therefore be chosen with each singleton norm at
most $T$. Scale to $S=3$. Equations~\eqref{eq:envelope-use}--\eqref{eq:Kbounds}
give, with $N_j=\Np(x_j)$,
\begin{equation}\label{eq:normal}
1\le T<A,\quad N_j\le T,\quad N_j>3-2A,\quad
G<\frac{2}{cp},\quad k_0G+T<3.
\end{equation}

In three real coordinates each pair has a signed coordinate at least
$3^{-1/p}$ times its norm. Since its norm is at most $2A$, its loss
is less than $e$ provided
\begin{equation}\label{eq:orientchecks}
2A(1-3^{-1/p_0})<e,\quad
3-2A>2/k_0+e,\quad 3-4/k_0-2e>A.
\end{equation}
These comparisons are certified. The three selected coordinates are
distinct. Indeed, if pairs $(a,b)$ and $(a,d)$ select one coordinate with
the same sign, adding their inequalities and using the singleton bound
for $a$ gives a signed total coordinate at least $3-2G-2e>T$.
If they select opposite signs, each inequality, after bounding the other
singleton coordinate, gives $\pm a_i\ge N_a-G-e$. Hence
$N_a\le G+e<3-2A$, again impossible.

Permute and flip coordinates so that the pair complementary to index $i$
has its positive main coordinate at $i$. Write $X_{ji}$ for coordinate
$i$ of vector $j$, $P_i=\Np(X_j+X_k)$ and $L_i=X_{ji}+X_{ki}$, where
$\{i,j,k\}=\{1,2,3\}$. The following rational linear system contains
the resulting counterexample:
\begin{equation}\label{eq:LP}
\begin{gathered}
0\le N_j\le T,\quad\sum_jN_j=3,\quad1\le T\le A,\\
-N_j\le X_{ji}\le N_j,\quad -T\le\sum_jX_{ji}\le T,\\
0\le P_i\le N_j+N_k,\quad P_i-L_i\le e,\quad
T-k_0\sum_iP_i\le3-6k_0.
\end{gathered}
\end{equation}
Eighteen rational dual witnesses imply $0\le L_i\le3$ and
$|(X_j+X_k)_r|\le L_i/2$ for $r\ne i$. Concavity of $s^{1/p}$ yields
\[
P_i\le L_i(1+2\,2^{-p})^{1/p}\le L_i+6\,2^{-p}/p<L_i+\varepsilon_0.
\]
The last bound is checked at $p_0$. Add these three linear inequalities
to \eqref{eq:LP}. Eighteen further dual witnesses imply the coarse box
\begin{equation}\label{eq:box0}
-u_0\le X_{ii}\le-\ell_0,\qquad v_0\le X_{ji}\le w_0\quad(j\ne i).
\end{equation}
For clarity, a dual witness for $Av\le b$, $Ev=f$ consists of rational
$y\le0,z$ with $A^Ty+E^Tz=d$. It proves $d^Tv\ge b^Ty+f^Tz$.
The verifier reconstructs all objectives and checks these identities
exactly; an optimizer's claimed optimum is never used as evidence.

\subsection{A smaller tail bound}
In any box with parameters $(\ell,u,v,w)$ as in \eqref{eq:box0}, a
non-main coordinate of the pair excluding $k$ equals $V_i-X_{ki}$,
where $V=\sum_jX_j$ and $i\ne k$. Since $X_{ki}\ge v$ and
$V_i\le\Np(V)\le A$, it lies in
\begin{equation}\label{eq:clip}
[v-u,\ \min(w-\ell,A-v)].
\end{equation}
Window I uses this clipped interval. Its coarse tail bound is $.516$,
and its main coordinate is at least $1.396$, giving ratio $129/349$.
Window II uses the box-only tail bound $.706$ and main bound $1.5$.
In both windows the certificate checks, with the corresponding ratio $r_0$,
\begin{equation}\label{eq:tail}
\frac6{p_0}r_0^{p_0}<\varepsilon,\qquad 0<r_0<1.
\end{equation}
Thus $P_i\le L_i+\varepsilon$ for all normalized counterexamples.
The clipping argument uses the known total norm cap; it makes no claim
about arbitrary points of the coarse box without that cap.

\section{Certified contraction}\label{sec:contract}
Let $a_i=-X_{ii}$, $w_i=\sum_{r\ne i}X_{ir}$, $W=\sum_iw_i=\sum_iL_i$,
and $\gamma_p=2\,2^{-1/p}$. H\"older gives
\begin{equation}\label{eq:row}
a_i^p+(w_i/\gamma_p)^p\le N_i^p,\qquad w_i\le\gamma_pN_i.
\end{equation}
From $P=6-G$ and the refined tails,
$W>6-2/(cp)-3\varepsilon$.
On each rational slab $[\alpha,\beta]$ choose certified rational bounds
\begin{align}
t_-&\le2^{-1/\alpha}\le t_\alpha^+,&2^{-1/\beta}&\le t_\beta^+,\nonumber\\
h&\ge2/(c\alpha)+3\varepsilon,&
D&\ge2/(c\alpha)+3\varepsilon-6(1-t_\alpha^+),\label{eq:budgets}\\
D&\ge0,&\delta&\ge D/(2t_-),\quad\delta\ge0.\nonumber
\end{align}
The comparison $2/c>6\log2$ is certified. It makes
$6\,2^{-1/p}+2/(cp)$ decreasing. Therefore the nonnegative residuals
$r_i=\gamma_pN_i-w_i$ satisfy
\begin{equation}\label{eq:residual}
\sum_i r_i<D,\qquad W>6-h,\qquad 2t_-\le\gamma_p\le2t_\beta^+.
\end{equation}
In particular, a pair of rows consumes the same total budget $D$,
not two separate copies of that budget.

\begin{lemma}[Diagonal and ratio bounds]\label{lem:contract}
Suppose $\delta\le N\le V$, $a^p+(w/\gamma_p)^p\le N^p$,
and $w\ge\gamma_pN-D$. If $B\ge0$ and
\begin{equation}\label{eq:diagcert}
(B/V)^p+(1-\delta/V)^p\ge1,
\end{equation}
then $a\le B$. Also, if $N_i\ge n>0$ and
\begin{equation}\label{eq:ratiocert}
\begin{gathered}
1\le s\le2t_--D/n,\qquad0\le\tau\le s/2\le\sigma,\\
\tau^p+(s-\tau)^p>1,\qquad\sigma^p+(s-\sigma)^p>1,
\end{gathered}
\end{equation}
with all four power bases in $[0,1]$, then each off-diagonal entry of
row $i$ lies strictly between $\tau N_i$ and $\sigma N_i$.
\end{lemma}
\begin{proof}
Since $D/\gamma_p\le\delta$, the first hypotheses imply
$a^p\le N^p-(N-\delta)^p\le V^p-(V-\delta)^p\le B^p$;
the middle function is increasing for $N\ge\delta$.
For the second claim divide the two off-diagonal entries by $N_i$,
obtaining $u,v\ge0$, $u^p+v^p\le1$, $u+v\ge s$. Both are at most one.
If $u\le\tau$, use the decrease of $r^p+(s-r)^p$ on $[0,s/2]$.
If $u\ge\sigma$, use its increase on $[s/2,s]$; $u\le1\le s$.
Either case contradicts the stated power comparison. The same applies to $v$.
\end{proof}

Initialize $n=v_0$, which is a valid lower bound on each $N_i$.
The following updates are applied sequentially, always using the already
established lower bound and checking $\delta\le n$.

\paragraph{Lower update.}
Choose $n'>0$ and $B\ge0$ satisfying \eqref{eq:diagcert} at $V=n'$ and
$A\le\tau(3-n')-B$. If $N_i<n'$, the diagonal bound gives $a_i\le B$,
whereas the $i$th total coordinate exceeds
$\tau(3-N_i)-B>A$, a contradiction. Hence $N_i\ge n'$.

\paragraph{Upper update.}
Choose $0<U<3$ and $B\ge0$ satisfying \eqref{eq:diagcert} at $V=3-U$ and
\begin{equation}\label{eq:upperupdate}
2t_\beta^+A\le6-h-\sigma(3-U)-B.
\end{equation}
If $N_i>U$, let $M=N_j+N_k<3-U$. Minkowski applied to
\eqref{eq:row} gives
\[
(a_j+a_k)^p+((w_j+w_k)/\gamma_p)^p\le M^p.
\]
The sum of their residuals is less than $D$, and $M\ge2n\ge\delta$.
The diagonal bound thus gives $a_j+a_k\le B$. The sum of total
coordinates $j,k$ exceeds
$W-\sigma M-B>6-h-\sigma(3-U)-B$.
H\"older bounds this sum by $\gamma_pT<2t_\beta^+A$, contradicting
\eqref{eq:upperupdate}. Thus $N_i\le U$.

\paragraph{Final box.}
With the final $n,U,\tau,\sigma$, choose $A_f$ satisfying the diagonal
certificate at $V=U$. Then
\begin{equation}\label{eq:finalbounds}
\tau n<X_{ij}<\sigma U\ (i\ne j),\qquad
\tau(3-U)-A<a_i\le A_f.
\end{equation}
Each slab stores a rational box $B=B(\ell,u,v,w)$, verified by
\begin{equation}\label{eq:boxchecks}
v\le\tau n,\quad \sigma U\le w,\quad A_f\le u,\quad
\ell\le\tau(3-U)-A.
\end{equation}
There are ten lower updates followed by the final ratio and upper update
in every slab. All power comparisons are checked at $p=\beta$ with bases
in $[0,1]$, so they hold throughout the slab. These are finite induction
steps, with no assumption of the desired final box in an earlier update.

\section{Convex domains retaining the normalization information}\label{sec:domain}
The normalization $S=3$ is useful for deriving bounds but is not itself a
convex constraint to preserve during averaging. Window II uses
\begin{equation}\label{eq:D2}
\mathcal D_2=B\cap\{\Np(X_i)\le U\ (i=1,2,3),\ \Np(V)\le A\}.
\end{equation}
For Window I, each pair gap is nonnegative and at most $G$; hence at a
normalized counterexample
\[
P_i=3-N_i-(N_j+N_k-P_i)>3-U-2/(c\alpha).
\]
By the tail bound, set
\begin{equation}\label{eq:floors}
J=3-U-2/(c\alpha)-\varepsilon,\qquad W_0=6-h.
\end{equation}
Then the counterexample belongs to the smaller domain
\begin{equation}\label{eq:D1}
\mathcal D_1=\mathcal D_2\cap\{L_i\ge J\ (i=1,2,3),\ \sum_iL_i\ge W_0\}.
\end{equation}
Both domains are convex and invariant under simultaneous permutation of
rows and coordinates. The new constraints in \eqref{eq:D1} are linear.
In particular, neither domain retains the equality $S=3$.

Here are the constants used to bound the Hessian. Set $m=\min(\ell,v)$,
$M=\max(u,w)$. In Window I set
\begin{equation}\label{eq:geo1}
\begin{gathered}
L=\max(2v,J),\quad s=L-u,\quad t=\min(2w-\ell,A),\\
\eta=\max\{|v-u|,|\min(w-\ell,A-v)|\},\quad b_0=\max(v,W_0/6).
\end{gathered}
\end{equation}
The interval in \eqref{eq:clip} is also checked to be nonempty.
In Window II use the weaker box bounds
\begin{equation}\label{eq:geo2}
L=2v,\quad s=2v-u,\quad t=2w-\ell,\quad
\eta=\max(|v-u|,|w-\ell|),\quad b_0=v.
\end{equation}
The certificate verifies $m>0$, $0<s\le t$, $0<\eta\le L$ and,
in Window I, $J,W_0>0$. On the appropriate domain, each singleton
coordinate has absolute value in $[m,M]$, total coordinates lie in $[s,t]$,
and each pair has main coordinate at least $L$ and tails of magnitude
at most $\eta$. Singletons, pairs, and total are therefore nonzero.

\section{The Hessian estimate}\label{sec:hessian}
Write $|\cdot|_2$ for the Euclidean norm. For $p>2$, $\Np$ is $C^2$
away from the zero vector, even if some individual coordinates vanish.
Differentiation and completion of the square give
\begin{equation}\label{eq:hessian}
H_p(v,z)=\frac{p-1}{\Np(v)^{p-1}}
\min_{a\in\R}\sum_i|v_i|^{p-2}(z_i-av_i)^2.
\end{equation}
The minimizing $a$ is
$\sum_i|v_i|^{p-2}v_iz_i/\sum_i|v_i|^p$.
For each row decompose a direction as $Z_j=a_jX_j+U_j$ using this
minimizer, and put $D_Z=\sum_j Z_j-bV$ using the total minimizer.
The norm caps in the convex domain imply the positive lower bounds
\begin{equation}\label{eq:positivehess}
\sum_jH_p(X_j,Z_j)\ge b_C\sum_j|U_j|_2^2,\quad
H_p(V,\sum_jZ_j)\ge b_T|D_Z|_2^2,
\end{equation}
where
\[
b_C=\frac{(p-1)m^{p-2}}{U^{p-1}},\qquad
b_T=\frac{(p-1)s^{p-2}}{A^{p-1}},\qquad
d=\frac{(p-1)\eta^{p-2}}{L^{p-1}}.
\]
In a pair with main coordinate $k$, choose in \eqref{eq:hessian} the
scalar which eliminates the main coordinate. The resulting upper bound
is $d$ times the sum of squares of its two remaining components.

\subsection{Controlling the radial coefficients}
Let $a_j=\max a_i$, $a_k=\min a_i$, $\Delta=a_j-a_k$, and
$R=\sum_i a_iX_i$. Then
\begin{equation}\label{eq:cross}
V_jR_k-V_kR_j\ge m(V_j+V_k)\Delta.
\end{equation}
Indeed, subtract $a_k$ from all coefficients, and let $l$ be the third
index and $\lambda=a_l-a_k\in[0,\Delta]$. The left side becomes
$\Delta A_1+\lambda A_2$, where
$A_1=V_jX_{jk}-V_kX_{jj}$ and $A_2=V_jX_{lk}-V_kX_{lj}$.
If $A_2\ge0$ it is at least $\Delta A_1$; otherwise it is at least
$\Delta(A_1+A_2)$. Both $A_1$ and
$A_1+A_2=V_kX_{kj}-V_jX_{kk}$ are at least $m(V_j+V_k)$.

Define
\begin{equation}\label{eq:Hshape} H=\frac{s^2+t^2}{m^2(s+t)^2}.\end{equation}
For two positive numbers in $[s,t]$, their squared Euclidean norm divided
by their squared sum is at most $(s^2+t^2)/(s+t)^2$; this follows by
monotonicity in their ratio. Since $R=D_Z-\sum U_i+bV$, Cauchy and
Young's inequality in \eqref{eq:cross} give
\begin{equation}\label{eq:delta}
\Delta^2\le H\left((1+e_1)|D_Z|_2^2+
3(1+e_1^{-1})\sum_i|U_i|_2^2\right).
\end{equation}

\subsection{Pair residuals}
For the pair of rows $j,r$ complementary to $k$, the residual in coordinate
$i\ne k$ equals $r_{ki}+e_{ki}$, where
\begin{align}
r_{ki}&=(a_j-a_r)\frac{X_{rk}X_{ji}-X_{jk}X_{ri}}{X_{jk}+X_{rk}},\nonumber\\
e_{ki}&=U_{ji}+U_{ri}-\theta_{ki}(U_{jk}+U_{rk}),\qquad
\theta_{ki}=\frac{X_{ji}+X_{ri}}{X_{jk}+X_{rk}}.\label{eq:pairres}
\end{align}
The fraction multiplying $a_j-a_r$ is bounded by $M$ in absolute value:
it is a difference with positive weights summing to one, and both relevant
coordinates have absolute value at most $M$. Thus
\begin{equation}\label{eq:rbound}
\sum_{k,i\ne k}r_{ki}^2\le2M^2\sum_{j<r}(a_j-a_r)^2\le4M^2\Delta^2.
\end{equation}
For Window II, $|\theta_{ki}|\le1$ and four-term Cauchy gives
$\sum e_{ki}^2\le16\sum_j|U_j|_2^2$: in the six residuals an off-diagonal
entry occurs at most three times and a diagonal entry twice.

For Window I put $\rho=\eta/L$. Split $e_{ki}$ into
$U_{ji}+U_{ri}$ and $-\theta_{ki}(U_{jk}+U_{rk})$.
In each coordinate, the first set contributes $(x+y)^2+(x+z)^2\le
3(x^2+y^2+z^2)$, since the difference is
$(x-y-z)^2+(y-z)^2$. The second set has squared norm at most
$4\rho^2\sum|U_j|_2^2$, because each off-diagonal main entry occurs
twice. The Euclidean triangle inequality yields
\begin{equation}\label{eq:ebound}
\sum e_{ki}^2\le(\sqrt3+2\rho)^2\sum_j|U_j|_2^2
\le E\sum_j|U_j|_2^2,\qquad E=(7/4+2\rho)^2.
\end{equation}
For Window II set $E=16$. Using Young's inequality with $e_2$,
\eqref{eq:delta} and \eqref{eq:rbound} give
\begin{align}
\sum(r_{ki}+e_{ki})^2&\le G_T|D_Z|_2^2+G_C\sum_j|U_j|_2^2,\label{eq:pairbound}\\
G_T&=4M^2(1+e_2)H(1+e_1),\nonumber\\
G_C&=12M^2(1+e_2)H(1+e_1^{-1})+E(1+e_2^{-1}).\nonumber
\end{align}
Combining the positive and negative Hessian bounds proves
\begin{equation}\label{eq:Fhessian}
D^2\F(X)[Z,Z]\ge(b_T-KdG_T)|D_Z|_2^2+
((2K-1)b_C-KdG_C)\sum_j|U_j|_2^2.
\end{equation}
Put
\begin{equation}\label{eq:ratios}
r_T=\frac A L\left(\frac{\eta A}{Ls}\right)^{p-2}=\frac d{b_T},\qquad
r_C=\frac U L\left(\frac{\eta U}{Lm}\right)^{p-2}=\frac d{b_C}.
\end{equation}
On every slab the certificate checks the two strict inequalities
\begin{equation}\label{eq:curvature}
1-G_T(p/2)r_T>0,\qquad 2c-1/p-(G_C/2)r_C>0.
\end{equation}
By \eqref{eq:Kbounds}, both coefficients in \eqref{eq:Fhessian} are
positive. In particular $\F$ is convex on the relevant domain: along
each segment it is $C^2$ in a neighborhood of every point and has
nonnegative second derivative. This also applies to segments on the
boundary of the closed domain.

\section{Permutation averaging and sharpness}\label{sec:finish}
Average $X$ over all six simultaneous row and coordinate permutations.
The resulting matrix has diagonal entries $-a$ and off-diagonal entries
$b>0$. It remains in the convex domain, and convexity gives
$\F(\overline X)\le\F(X)$. The box gives $a\le u$, $b\ge v$;
Window I also gives $b\ge W_0/6$, hence $b\ge b_0$ in both windows.
The certificate verifies
\begin{equation}\label{eq:orbit}
u/v\le2,\qquad \ell/w\ge1/2\quad\hbox{or}\quad A\le3b_0/2.
\end{equation}
If the box alternative holds, $a/b\ge1/2$ directly. Otherwise the
positive total coordinates give $0<2b-a\le A$, and therefore
$a/b\ge2-A/b\ge1/2$. Thus $t=a/b\in[1/2,2]$.
For $K=K_p$, equation~\eqref{eq:cyclic} implies
\[
\F(\overline X)=b(6A_p(t)-3B_p(t))(K_p-R_p(t))\ge0.
\]
A counterexample would instead give $\F(\overline X)\le\F(X)<0$.
This excludes real three-dimensional counterexamples on both windows.
Lemma~\ref{lem:dim} proves the complex inequality on $[21,32]$.
The accepted theorem \cite{solution32} supplies $p\ge32$ with the same
definitions, proving admissibility in Theorem~\ref{thm:main}.

Finally, the cyclic triple has strictly positive $G$ and ratio exactly
$R_p(t)$. A constant valid in complex dimension three must dominate every
such ratio, and hence $K_p$. Taking a maximizer shows attainment. Padding
with zeros gives the same sharpness in each dimension $n\ge3$.

\section{Finite certificates and reproducibility}\label{sec:cert}
\begin{lemma}[Finite certificate lemma]\label{lem:finite}
For the two profiles in Table~\ref{tab:profiles}, all scalar, orientation,
linear-dual, tail, contraction, geometry, orbit, and curvature comparisons
used in Sections~\ref{sec:coarse}--\ref{sec:finish} hold on their specified
closed real exponent intervals.
\end{lemma}
\begin{proof}[Computer-assisted verification]
The proof is the deterministic replay described here, using the committed
JSON witnesses and the directed interval operations of Appendix~\ref{app:arith}.
Table~\ref{tab:counts} records the finite size. All profile parameters,
partition endpoints, and witnesses are parsed as exact fractions.
Partitions must meet exactly, in order, with positive length and matching
outer endpoints. Thus all real exponents are covered.

For each scalar cell $[a,b]$, the witness supplies one rational
$t\in[1/2,2]$. Evaluate $R_p(t)-E_p(q)$ and $R_p(t)-cp$ with
$p\in[a,b]$ and require positive interval lower endpoints. This proves
the lower bounds in \eqref{eq:Kbounds}. For the upper bound, partition
the exponent range and, on each exponent cell, partition the entire
$t$ range. Evaluate $p/2-R_p(t)$ on every rectangle and require a positive
lower endpoint. Every denominator is checked positive before division.

The coarse checker reconstructs \eqref{eq:LP} and the 36 objectives;
it checks each rational dual identity, sign and target, adding the initial
tail constraints only between its two stages. The scalar orientation,
tail and monotonicity prerequisites are checked independently.
For each contraction slab it validates \eqref{eq:budgets}, the preceding
$n$ in each round, the bases and branch conditions in
Lemma~\ref{lem:contract}, the updates, and \eqref{eq:boxchecks}.
It then reconstructs all geometric constants as rational numbers and
checks their domains and \eqref{eq:orbit}. Finally, interval evaluation
of \eqref{eq:ratios}--\eqref{eq:curvature} covers the full slab.
Stored floating-point diagnostics are not used in any acceptance decision.
\end{proof}

\begin{table}[ht]
\centering
\begin{tabular}{lrr}\toprule
Obligation & $[21,26]$ & $[26,32]$\\\midrule
Scalar exponent cells &2355&2699\\
Upper-bound rectangles &960&672\\
Exact LP dual witnesses &36&36\\
Contraction slabs &160&96\\
Lower updates per slab &10&10\\\bottomrule
\end{tabular}
\caption{Certificates cover intervals, not isolated sample points.}\label{tab:counts}
\end{table}

The primary files are \path{profiles/cutoff21.json},
\path{profiles/cutoff26.json} and the corresponding subdirectories of
\path{certificates/}, each containing \path{scalar.json},
\path{cyclic-upper.json}, \path{coarse.json}, and \path{contractor.json}.
The acceptance implementation is \path{scripts/verify.py}, with exact
model construction in \path{scripts/coarse_lp.py} and geometric constants
in \path{scripts/geometry.py}. Candidate generation may use floating-point
optimization, but the replay neither calls that optimizer nor trusts its
reported optimum.

Both windows have passed Arb at 192 and 256 bits and the standard-library
rational backend at 256 bits. The minimum reported lower endpoints of
the two dimensionless curvature margins on $[21,26]$ are approximately
$0.391298$ and $0.321741$; on $[26,32]$ they exceed $0.3200$ and $0.2675$.
These rounded summaries are descriptive; the actual decisions use the
unrounded interval endpoints.

From the repository root, the complete dependency-free release replay is
\begin{verbatim}
python -S scripts/replay_release.py
\end{verbatim}
It runs both windows at 256-bit rational precision, rejection controls,
and checks all eleven preserved upstream file hashes. It writes
\path{results/release-replay.json}, including the input hashes. Python 3.14
was used for the recorded runs. No third-party package is imported by
this replay; \texttt{-S} disables site initialization.
For an individual window, for example:
\begin{verbatim}
python -S scripts/verify.py --backend rational --precision 256 \
  --profile profiles/cutoff21.json \
  --directory certificates/cutoff21 \
  --output results/cutoff21-rational.json
\end{verbatim}
The backslash continuation above is for a POSIX shell; on Windows enter
the command on one line. An Arb replay uses the same command without
\texttt{-S}, with \texttt{--backend arb}, after installing
\texttt{python-flint==0.9.0} (see \path{requirements.txt}).

Negative controls deliberately damage coverage, a cyclic witness, a
contraction chain, an LP dual, a certificate set, a final box, and a
curvature configuration. Window I additionally tests an unjustifiably
small clipped-tail budget. Both backends reject all eight Window I
mutations; Window II has seven. There are also 53 arithmetic cross-checks
against higher-precision Arb. Rejection tests are implementation checks,
not substitutes for the analytic lemmas.

\subsection{Trust and provenance}
The two arithmetic backends share the certificate data, LP model,
geometric formulas, and \path{verify.py}. They are not independent
end-to-end proofs. Python's integer arithmetic, interpreter, the interval
implementation, and the correspondence between this paper and the checker
remain in the computational trust boundary. The mathematical argument
has received internal self-audit, but no independent analytic review is
claimed. No new Lean proof or Prove2Me submission has been made.

The tail reference is the final \texttt{theorem solution} in submission
\path{2d248eb1-4938-4dee-8392-df24d2c70f2f}, with statement
``for every real $p\ge32$, the cyclic constant is the least uniform
complex $\ell_p$ Hlawka constant.'' Its record and byte-preserved source
are in \path{sources/}; \path{sources/provenance.json} gives identifiers
and SHA-256 digests. The coarse, contraction and convexity components at
cutoff $32$ are attributed in those records to \texttt{@bwdaugherty}.
The underlying public construction repository is \cite{savarin}, with
reference commit \path{1829591140522193a81ae721330e4d775c81123e}.
We do not count a source download or acceptance record as a local rebuild.

\section{Discussion}
The number $21$ is a sufficient lower threshold for the exponent $p$,
not the value of $K_p$ and not a proven optimal threshold. The constant
itself remains the exact real maximum \eqref{eq:cyclic}; a decimal
approximation to it would require a specified exponent and precision.
The proof does not require $p$ or the cutoff to be an integer. The older
window $[127/4,32]$ illustrates that rational cutoffs work in the same
certificate framework.

The gain comes from keeping information during convexification:
singleton and total norm caps reduce Hessian denominators; the total norm
clips pair tails; pair-main-coordinate floors improve the radial scale;
and the off-diagonal sum floor keeps the permutation average in the
cyclic parameter range. These constraints are convex or linear, which is
essential for Jensen averaging. Further threshold reductions would
require new certificates and, if estimates change, a corresponding
analytic justification. No claim for $2<p<21$ follows from the present
calculation.

\appendix
\section{Directed dyadic arithmetic}\label{app:arith}
At precision $B$ an integer pair $(l,u)$ represents
$[l2^{-B},u2^{-B}]$. Fraction inputs are enclosed using integer floor and
ceiling. Addition and subtraction use endpoint identities. Multiplication
and division use all endpoint combinations and outward integer rounding;
division rejects intervals containing zero. Float inputs are rejected.
Strict positivity is accepted only when the integer lower endpoint is
strictly positive. These rules give containment by elementary interval
arithmetic, regardless of dependency-induced overestimation.

For a positive dyadic $x$, range reduction writes $x=2^km$, with
$1\le m<2$. Let $z=(m-1)/(m+1)$; the implementation checks a rounded
enclosure in $[0,1/2]$. The finite formula
\[
\log m=2\sum_{j=0}^{191}\frac{z^{2j+1}}{2j+1}+R,\qquad
0\le R\le\frac{2z^{385}}{385(1-z^2)}
\le\frac8{3\cdot385\cdot2^{385}}
\]
follows by integrating the geometric series or the power series for
$\operatorname{arctanh}z$. All terms are evaluated outward, and the
rational remainder is added to the upper endpoint. Compute $\log2$ by
the same formula and use $\log x=k\log2+\log m$, with signed interval
multiplication when $k<0$.

For $e^x$, choose $k\ge0$ so that the interval for $z=x/2^k$ lies
in $[-1/2,1/2]$, and evaluate
\[
e^z=\sum_{j=0}^{80}\frac{z^j}{j!}+R,\qquad
|R|\le\frac2{2^{81}81!}.
\]
The Lagrange remainder and $e^{1/2}<2$ justify this bound.
Enlarge both endpoints by the remainder bound, verify positivity,
and square $k$ times outward. Monotonicity gives interval exponential
and logarithm from their endpoint enclosures. For positive bases use
$x^y=\exp(y\log x)$. When the base interval includes zero and its upper
endpoint is positive, with strictly positive exponents, monotonicity in
the base gives lower endpoint zero and the upper-base power enclosure.
A zero base gives zero. Other nonpositive-base cases are rejected.
This includes the vanishing factor $|1-t|^p$ at $t=1$.
The implementation is \path{scripts/rational_intervals.py}; explicit
remainders, not presumed numerical precision, supply rigor.

\section*{Acknowledgments and AI-use disclosure}
The cyclic reduction and formalized baseline are credited to the
\texttt{savarin/hlawka-schatten} project (Ezzeri Esa / savarin in the
source attribution) and the Prove2Me contributors, including
\texttt{@bwdaugherty}. Their sources and licenses are identified in
\path{sources/NOTICE.md}. This project adapts that proof architecture;
it does not claim to originate the diagonal Hlawka problem or the
cyclic extremizer construction.

OpenAI Codex assisted with mathematical exploration, proof drafting,
certificate generation, implementation, verification runs, and preparation
of this manuscript. AI assistance and internal replay do not constitute
independent expert review. The author name and institutional affiliation
are supplied by the author; no institutional endorsement is implied.

\section*{Data and software availability}
The manuscript source, PDF, exact input certificates, replay scripts,
and verification reports are available at \repo.
This is a repository preprint, not a journal publication. Citation
metadata and third-party notices accompany the release.

\begin{thebibliography}{9}
\raggedright
\bibitem{AK} K.~M.~R. Audenaert and F. Kittaneh,
\emph{Problems and Conjectures in Matrix and Operator Inequalities},
arXiv:1201.5232v3 (2012), Section 8.2, Problem 7.
\url{https://arxiv.org/abs/1201.5232v3}.
\bibitem{MN} D.-\c{S}. Marinescu and C.~P. Niculescu,
\emph{A survey of the Hornich--Hlawka inequality},
arXiv:2407.03278 (2024). \url{https://arxiv.org/abs/2407.03278}.
\bibitem{savarin} savarin and contributors,
\emph{Dimension-independent Hlawka constants for Schatten norms},
Lean source repository (2026). \url{https://github.com/savarin/hlawka-schatten}.
Reference commit recorded in the accompanying provenance file.
\bibitem{campaign} Prove2Me contributors,
\emph{Sharp diagonal Hlawka constant}, campaign record, accessed
17 October 2026. \url{https://prove2.me/campaigns/sharp-diagonal-hlawka-constant}.
\bibitem{solution32} Prove2Me contributors,
\emph{Accepted solution for cutoff 32}, submission
\path{2d248eb1-4938-4dee-8392-df24d2c70f2f} (16 October 2026),
theorem record \url{https://prove2.me/theorems/6e8430ce-7c32-4b82-9a1b-d7fcd67fdeb3}.
Source and acceptance metadata preserved under \path{sources/}.
\end{thebibliography}
\end{document}
